\documentclass[../report.tex]{subfiles}

\begin{document}

\subsection{Windows Scanning and Buffer Overflow} \label{buff:vulns}

\subsubsection{Vulnerability - Buffer Overflow} \label{subsubsec:buffer_overflow_expl}
To access the system initially, a Buffer Overflow vulnerability was used (see \fullref{subsec:buffer_overflow} for execution details).
The vulnerability can be classified as:\\
\textbf{CWE-120: Buffer Copy without Checking Size of Input ('Classic Buffer Overflow')} \autocite{mitre:cwe}\\
There is no CVE entry for the vulnerability in the VulnServer program, since it was specifically designed as an educational tool for Buffer Overflows \autocite{vulnserver}.

\paragraph{Cause and Prevention}
The VulnServer application from Bradshaw's Github repository uses \emph{strcpy} (see \cref{strcopy}), a C function which does not check if enough space is present before copying input to the buffer \autocite{microsoft:strcpy}. 

\begin{figure}[!h]
	\centering
	\includegraphics[width=0.6\linewidth]{buffer_overflow/analysis/strcopy}
	\caption{Example of C Code vulnerable to Buffer Overflows from VulnServer Repo}
	\vspace{-2mm}
	\caption*{\small Source: \url{https://github.com/stephenbradshaw/vulnserver/blob/master/vulnserver.c}}
	\label{strcopy}
\end{figure}

The modified version of VulnServer deployed in the laboratory similarly does not conduct any bounds checks when processing the input, allowing for the Buffer Overflow to occur.
To fix this vulnerability, a safer function like strcpy\_s could be used \autocite{microsoft:strcpy}, or a bounds check could be implemented before calling the unsafe function. In the latter case, the input length should be validated, ensuring it fits inside the allotted char buffer without flowing over.

\subsubsection{Risk Assessment}

%\begin{figure}[!h]
%	\centering
%	\includegraphics[width=0.75\linewidth]{buffer_overflow/analysis/cvss}
%	\caption{CVSS v4.0 Score calculation}
%	\vspace{-2mm}
%	\caption*{\small Source: \url{https://www.first.org/cvss/calculator/4.0}}
%\end{figure}

\begin{tabular}{l l}
\textbf{Vector}: & CVSS:4.0/AV:N/AC:L/AT:N/PR:N/UI:N/VC:N/VI:N/VA:H/SC:H/SI:H/SA:H\\
\textbf{CVSS v4.0 Score}: & 9.3 / Critical\\
\end{tabular}

\paragraph*{Exploitability Metrics}
The vulnerability works via \textit{Network (N)} by connecting to the server, has a \textit{Low (L)} \textit{Complexity (C)} since execution of the script is all that is required, necessitates no \textit{Attack Requirements (AT)} other than the server being up, requires no \textit{Privileges (P)}, and no \textit{User Interaction (UI)} to work.

\paragraph*{Impact Metrics}
The vulnerable system (VulnServer) can be crashed by the attack, scoring \textit{High (H)} for \textit{Availability (VA)}. \textit{Confidentiality (VC)} and \textit{Integrity (VI)} are not endagered in the same manner, since attack the server itself does not divulge any sensitive data to the attacker, nor do  files change because of the Buffer Overflow.

Since compromise of the vulnerable system (VulnServer) lead to a reverse shell on the subsequent system (Windows 7 machine), \textit{Confidentiality (SC)}, \textit{Integrity (SI)}, and \textit{Availability (SA)} should be considered highly endangered on the subsequent system. The privileges the actual VulnServer process would run with in a real world scenario are not know, but it is safe to assume that reverse shell access to the system is a highly threatening to the CIA triad.

\newpage
\subsubsection{Countermeasures} \label{counta}
\begin{itemize}
\itemBF{Exploit Protection (M1050)}
Necessitating DEP and enabling ASLR for the VulnServer process would have prevented the abuse of the permanent JMP ESP location command, making repeated exploitation more difficult.
\itemBF{Filter Network Traffic (M1037)}
Monitoring network traffic can help identify threats and prevent them.
IP addresses could be blacklisted via firewall e.g. if port scanning is detected (\textit{DET0830, DET0867}).
Monitoring outbound callbacks could also be of value. A webserver shouldn't connect to port 443 on another machine, that could be blocked by a firewall rule  (\textit{DET0080, M1037, DET0444}).
\itemBF{Application Isolation and Sandboxing (M1048)}
If an attacker were to exploit a vulnerability, sandboxing the application could make it more difficult for the whole system to get compromised.
\itemBF{Network segmentation (M1030)}
Create a demilitarized zone (DMZ) to encapsulate public-facing services can limit damages to the whole network if an attack were to be successful.
\itemBF{Exploit Protection (M1050)}
Employing a SIEM system or other security applications (like IDS / IPS) could aid detection of malicious activity.
\itemBF{Privileged Account Management (M1026)}
Use the least privilege principle to limit exploitable permissions.
\itemBF{Vulnerability Scanning (M1016)}
Periodically scan public-facing systems for vulnerabilities. This could lead to the discovery of potential weaknesses like the Buffer Overflow.
\end{itemize}

\end{document}